% ---
% id: reductions/wpfss/rem:published
% kind: rem
% title: "Relation to the published flag-state squashers"
% summary: >-
%   The flag-state target is the flag-state squasher of Zhang et al., and the
%   single-bound weight-preserving target that of Nahar et al. Their usual
%   scalar flag-weight constraint is a carried-operator entry stronger than the
%   default licensed one, and only its operator form is ever emitted.
% prominence: standard
% status: accepted
% uses: [def:targets, prop:chain, reductions/cor:transported, reductions/cor:pullback]
% mentions: [reductions/cor:entries, cor:two-sector-decided, cor:fssinstance, reductions/def:exactadm, reductions/req:declared-constants-only, reductions/req:licensed-entry-lattice]
% assumes: [A-CONST]
% refs:
%   - "zhang2021security, Thm. 1, Eq. (8)"
%   - "tupkary2026decoy, Eqs. (149)-(150)"
%   - "nahar2024postselection, Lemma 7, Eqs. (28)-(29)"
%   - "nahar2024postselection, Eq. (27)"
% provenance: newly-derived
% lean:
% ---

\begin{remark}
\label{rem:published}
$\Ffss$ is the flag-state squasher of \cite[Thm.~1, Eq.~(8)]{zhang2021security},
which states exactly \eqref{eq:FSS}: $\widetilde M_y=M_{y,n\le k}\oplus
\dyad{y}$ for a POVM block diagonal at a photon-number cutoff. The adjoint
form is stated there in part: \cite[Eqs.~(149)--(150)]{tupkary2026decoy}
gives, for that map, the defining identity $\Lambda^{\dagger}(F_i)=M_i$ and the
\emph{flag} half of \eqref{eq:transportproj},
$\Lambda^{\dagger}(\Pflag)=\Pgt$, and states no second projector identity. The
preserved half follows from it in one step and is written out here rather than
attributed: $\Lambda$ is trace preserving, so $\Lambda^{\dagger}$ is unital, and
$\Pleq+\Pflag=\idon{\Qsys}$ with $\Pleq+\Pgt=\idon{\Bsys}$ give
$\Lambda^{\dagger}(\Pleq)=\idon{\Bsys}-\Lambda^{\dagger}(\Pflag)
=\idon{\Bsys}-\Pgt=\Pleq$. This contract proves both halves at
Proposition~\ref{prop:chain}\,(i) and the citation is corroboration, which is
why the narrower attribution costs nothing. The \tref{single-bound case} of
the \tref{weight-preserving target} $\Fwp$
is the weight-preserving flag-state squasher of
\cite[Lemma~7, Eqs.~(28)--(29)]{nahar2024postselection}, whose target places
$f_N$ on the diagonal of the flag block of one distinguished element and
$(1-f_N)$ on the others; the multi-bound form \eqref{eq:WPFSS} is the
symmetric statement with one weight per outcome.

The scalar flag-weight constraint these targets are usually paired with is
\cite[Eq.~(27)]{nahar2024postselection}, and the comparison with what this
module emits runs in the direction opposite to the naive one. Writing
$a\coloneqq\lmin\big(\blk{M}{\leqN}_e\big)$ and
$b\coloneqq\lmin\big(\blk{M}{\gtN}_e\big)$ for a designated event $e$
(declared device constants, fixed before the run and never recomputed from
the declaration, Assumption~\ref{ass:constdecl}), that
equation reads, in full, $W\le1-p(m\le N)\le\big(p(e)-a\big)/(b-a)$ for the
flag weight $W=\tr[\Pflag\sigma]$ and the observed
$p(e)=\tr[\Ffss_e\sigma]$;
the right-hand member is the operative one and the only one compared with
below. It is
\emph{a} member of the family Corollary~\ref{cor:transported} generates,
namely the carried operator $\hat O$, which carries the certificate
\begin{equation}\label{eq:two-sector-entry}
\begin{split}
  \hat O&=\Ffss_e-a\,\Pleq-b\,\Pflag,\\
  \Lfss^{\dagger}\big(\hat O\big)&=M_e-a\,\Pleq-b\,\Pgt
  =\big(\blk{M}{\leqN}_e-a\idon{}\big)
    \dsum\big(\blk{M}{\gtN}_e-b\idon{}\big)
  \;\succeq\;0
\end{split}
\end{equation}
by the definitions of $a$ and $b$ and the transported projectors of
\eqref{eq:transportproj}. The emitted constant is $0$, with no shift
(Corollary~\ref{cor:entries}\,(iii)).
Corollary~\ref{cor:pullback}\,(ii) gives
$p(e)-a(1-W)-bW\ge0$, which rearranges to Eq.~(27) when $b>a$ (the published
division is by $b-a$; at $b\le a$ it reverses or is undefined, and only the
operator form above is used here). It is \emph{not} the member
Proposition~\ref{prop:chain}\,(iii) supplies and not the one
Corollary~\ref{cor:fssinstance} emits: those use $a=0$, where the same
corollary gives only $p(e)-bW\ge0$.

Which of the two is stronger is settled at the level of the \emph{constraint
sets}, not by comparing the two caps on $W$: since
$\hat O=\hat O\big|_{a=0}-a\,\Pleq$ with $a\ge0$ and $\Pleq\succeq0$, one has
$\tr[\hat O|_{a=0}\,\sigma]\ge\tr[\hat O\,\sigma]$ for every $\sigma\succeq0$,
so $\{\sigma:\tr[\hat O\sigma]\ge0\}\subseteq
\{\sigma:\tr[\hat O|_{a=0}\sigma]\ge0\}$: the $a>0$ entry cuts a set no
larger and is at least as strong a constraint, unconditionally. The implication
therefore runs from Eq.~(27) to the $a=0$ form and not the other way. (The
comparison of the two \emph{caps} $\big(p(e)-a\big)/(b-a)$ and $p(e)/b$ is not
the right one and does not run in a fixed direction: the former is the smaller
exactly when $p(e)<b$, because $p(e)$ is itself a function of $\sigma$ rather
than a fixed number, so the two inequalities cut the state set differently.)

Two consequences follow. \emph{(a)} The default licensed entry of
Corollary~\ref{cor:fssinstance} is weaker than standard practice, not
stronger; a declaration may request the entry \eqref{eq:two-sector-entry}
as a carried operator with declared constants $a$ and $b$, which need not be
the least eigenvalues above. Its certificate is then not taken from the
request: it is decided on the emitted numbers, where every emitted block of
$e$ below the cutoff clears $a$ and $b$ is at most the certified weight
$p_e$ (Corollary~\ref{cor:two-sector-decided}), and the request is refused
where either fails. Both constants must be declared on the lattice of
\eqref{eq:latinv}, as every declared constant must: an $a$ or $b$ off the
lattice puts the requested operator's entries off the lattice, and the
request is refused (Requirement~\ref{req:licensed-entry-lattice}). Smaller
lattice values of $a$ and $b$ keep both conditions. \emph{(b)} The \emph{cap} in
the displayed form is never emitted: its right-hand side contains the run's
observed $p(e)$, whereas Assumption~\ref{ass:constdecl} admits only constants
declared before the run, and an entry with such a constant is refused
(Requirement~\ref{req:declared-constants-only}). The admissible object is the
operator entry \eqref{eq:two-sector-entry} with $\bowtie\,=\,\succeq$ and $c=0$, of
which the cap is a run-time consequence and not a declaration. Every statement
of this remark is proved in the claims it names
(Definition~\ref{def:targets}, Proposition~\ref{prop:chain},
Corollary~\ref{cor:fssinstance}, Corollary~\ref{cor:two-sector-decided} and
Corollary~\ref{cor:entries}); the
citations record where a reader finds the correspondence.
\end{remark}
